\chapter{File formats and data}
\label{ch:formats}

\section{Products}

\begin{center}
\begin{tabularx}{\textwidth}{@{}lX@{}}
\toprule
File & Content \\
\midrule
\file{name.fbk.md} & the report: summary, analysis sections, findings,
  provenance, references (complete citations + paste-ready BibTeX) \\
\file{name.fbk.json} & the same content as data: \code{title}, \code{sections},
  \code{findings} (severity, message, rule id, suggested fix, refusals),
  \code{provenance}, plus a few program-specific keys \\
\file{structure.fbk.inp} & a generated ORCA input (pure ASCII) \\
\file{input.fix\_\textit{variant}.inp} & a corrected input from \menu{9} \\
\bottomrule
\end{tabularx}
\end{center}

Nothing carries timestamps; products are deterministic.

\section{Input JSON: cross-level records (menu 7)}

\begin{codeblock}
[
  {"label": "No.1",
   "occupation": "(5f)^1 ...",        # optional, informational
   "energies": {"HF": -1930.87692, "CCSD(T)": -1932.97951}},
  {"label": "No.2", "energies": {"HF": -1930.87392, "CCSD(T)": -1932.97745}}
]
\end{codeblock}

The top level may also be an object with a \code{records} list. All records
must carry the same level names; labels must be unique.

\section{Input JSON: crystal-field fit (menu 10)}

\begin{codeblock}
{
  "point_group": "C3",                  # C3 | Oh | C1/none
  "J": 8.0,                             # the |J M> manifold
  "levels": [0.0, 12.3, ...],           # one energy per sampled state
  "coefficients": [[[re, im], ...], ...],  # amplitudes on |J M>, M = m - J
  "declaration": {                      # optional: the five A5 items
    "convention": "...", "projection": "...", "units": "cm^-1",
    "z_axis": "...", "origin": "..."},
  "compare": {                          # optional: a second parameter set
    "label": "...",
    "parameters": {"2,0": 1879.0, "2,2": -43.0},
    "declaration": {"projection": "..."}}
}
\end{codeblock}

Parameter keys are \code{"k,q"} strings or two-element lists. Pass the
amplitudes themselves, not their complex conjugates (the conjugation trap,
Chapter~\ref{ch:menu10}).

\section{Input JSON: magnetic-doublet table (menu 16)}

\begin{codeblock}
{
  "system": "...",                       # optional label
  "reference": 0,                        # optional: index or label
  "doublets": [
    {"label": "ground", "g": [0.41, 0.44, 9.04], "energy": 0.0,
     "axis3": [0.0, 0.0, 1.0]},          # unit axis of the largest g value
    {"label": "1st excited", "g": [0.41, 0.44, 9.04], "theta3": 9.52}
  ]
}
\end{codeblock}

Per doublet, \code{theta3} (degrees) or \code{axis3} says where the largest
$g$ axis points; the reference doublet may omit both (its \code{theta3} is 0
by definition when the other angles are computed from axes, and it must carry
\code{axis3} then). \code{energy} is optional and printed when present.

\section{Structure and charge input (menus 2 and 11)}

Standard XYZ (atom count, comment, \code{element x y z}). For \menu{11}, the
interactive prompts take per-element charges (\code{O=-2,H=0.4}) and optional
radial moments \code{r2,r4,r6} in \AA$^k$. An element present in the structure
but missing from the charge list is refused, with the missing symbols named.

\section{Knowledge data (inside the package)}

\begin{center}
\begin{tabularx}{\textwidth}{@{}lX@{}}
\toprule
File & Content \\
\midrule
\file{knowledge/sources.bib} & the reference database (BibTeX). Every
  literature-evidence item must carry a bibkey that resolves here; loading
  fails otherwise \\
\file{knowledge/rules/*.yaml} & the rule base: id, kind (recipe/diagnosis),
  condition tree, action, refusals, evidence, confidence, severity (diagnosis
  only) \\
\file{knowledge/basis\_ecp.yaml} & the basis/ECP recommendation table \\
\file{knowledge/templates/*.j2} & the input templates (ORCA, OpenMolcas) \\
\file{knowledge/toolindex/tools.yaml} & the external-tool index (relation and
  status per entry) \\
\file{ui/menu.yaml} & the menu definition: number, title, handler (append-only
  numbers) \\
\bottomrule
\end{tabularx}
\end{center}

\section{The report as an interface}

The Markdown report is for humans; the JSON companion is the scripting
interface. Its stability follows the release policy: within a minor version
the section titles and finding fields are stable; new fields may be added.
Finding \code{rule\_id}s are stable identifiers you may key on.
